\scanpage{16}
Un cas intéressant est celui où $\deleted{\unsure{K|X}}\unsure{X|K}$ \deleted{\unsure{admet}} \inserted{est \unsure{le|la} sup \unsure{d'un|de un} \unsure{nb} fini de sous-}\deleted{un nb fini de \unsure{pts} maximaux, et}objets modèles, ce qui signifie que l'ens des éléments maximaux de $K\inserted{=\unsure{\widetilde X|X}}$ est fini, et que tt élément de $K$ est majoré par un élément maximal. Pour un ens. ordonné \deleted{ayant \unsure{ces}} \inserted{la $2^{\text{\unsure{ième|iem}}}$} \unsure{propriété|propriétés}, un $M$-épinglage est connu quand on connait les $u^K_{\deleted{\xi}\alpha}$ pour $\deleted{\xi}\alpha\inserted{\in K}$, maximal.

La condition pour que des donn.
\[
u^K_{\deleted{\xi}\alpha}:\ {\unsure{I|J}}_{\tau_{\unsure{n|K}}(\deleted{\xi}\alpha)}
   \xrightarrow{\sim} K_{\leq\deleted{\xi}\alpha}\qquad
   (\unsure{\xi}\in\operatorname{Max}K,
\]
définissent un \unsure{êpinglage|épinglage}, est que pour $\deleted{\unsure{\xi|\forall}}$ $\alpha,\beta\in\operatorname{Max}K$ et $\xi\in K$, $\xi\leq\alpha,\beta$, \deleted{posant} $n=\tau_K(\xi)$, \deleted{\unsure{des|les}} deux iso
\[
{\unsure{I|D}}_n\xrightarrow{\sim}K_{\leq\xi}
\]
définis par l'épinglage \deleted{\unsure{de|des}} $ {\deleted{\unsure{\tau}}u}^K_\alpha$ de $K_{\leq\alpha}$, et par l'épinglage $u^{\unsure{K|k}}_\beta$ de $K_{\leq\beta}$, \deleted{\unsure{soient}} les \unsure{mêmes}.
% archive sheet number: 15
